% GENERATED FILE. Edit canonical lessons, then run npm run generate:books.
% canonical-source-hash: fnv1a64:1880b6a500230e7a
% canonical-lessons: GE-C115-kamm, GE-C115-muelleimer, GE-C115-besen, GE-C115-waschmaschine, GE-C115-sessel

\chapter{Comb, Bin, Broom, Washing machine, Armchair}
\label{ch:ge-comb-bin-broom-washing-machine-armchair}
\hlchaptermodality{115}

\begin{chapteropening}
\textbf{By the end of this chapter:} \emph{I can say a comb, a bin, a broom, a washing machine and an armchair.}

Complete the last lesson of chapter 115: I can say a comb, a bin, a broom, a washing machine and an armchair.
\end{chapteropening}

\section[der Kamm]{der Kamm — a comb}

\label{lesson:GE-C115-kamm}



\begin{quote}
Before the new one: say the German for a toothbrush, then the German for toothpaste.
\end{quote}

\subsection*{You'll want to know: der Kamm}
\textbf{der Kamm} — \textquotedblleft{}a comb\textquotedblright{}.

\begin{cousinweb}
\textbf{Kamm} $\leftarrow$ Old High German \emph{chamb} $\leftarrow$ Proto-Germanic \emph{*kambaz}, \textquotedblleft{}comb\textquotedblright{} — inherited, the same word as English \emph{comb}. (EtymDB 2.1)
\end{cousinweb}

\begin{grammarlens}[title={What you've built}]
One more word for this chapter.
\end{grammarlens}

\subsection*{Your turn}
\begin{itemize}
\raggedright
  \item \emph{Say it:} \emph{der Kamm}
  \item \emph{Say it:} \emph{der Kamm}, once more
  \item \emph{Say it:} say \emph{der Kamm}
  \item \emph{Recall:} say the German for a toothbrush, then the German for toothpaste, then say \emph{der Kamm} again
\end{itemize}

\begin{tcolorbox}[breakable,colback=teal!4,colframe=teal!35!black,title={Before you move on}]
What does der Kamm mean? (\textquotedblleft{}A comb\textquotedblright{}.) Say it once more.
\end{tcolorbox}
\section[der Mülleimer]{der Mülleimer — a bin}

\label{lesson:GE-C115-muelleimer}



\begin{quote}
Before the new one: say the German for toothpaste, then the German for a comb.
\end{quote}

\subsection*{You'll want to know: der Mülleimer}
\textbf{der Mülleimer} — \textquotedblleft{}a bin\textquotedblright{}.

\begin{cousinweb}
\textbf{Mülleimer} is the word for \textquotedblleft{}rubbish\textquotedblright{} + the word for \textquotedblleft{}bucket\textquotedblright{} (EtymDB 2.1). The bucket word, Middle High German \emph{eimer}, goes back to Latin \emph{amphora}, reshaped as if it meant \textquotedblleft{}one-carrier\textquotedblright{}.
\end{cousinweb}

\begin{grammarlens}[title={What you've built}]
One more word for this chapter.
\end{grammarlens}

\subsection*{Your turn}
\begin{itemize}
\raggedright
  \item \emph{Say it:} \emph{der Mülleimer}
  \item \emph{Say it:} \emph{der Mülleimer}, once more
  \item \emph{Say it:} say \emph{der Mülleimer}
  \item \emph{Recall:} say the German for toothpaste, then the German for a comb, then say \emph{der Mülleimer} again
\end{itemize}

\begin{tcolorbox}[breakable,colback=teal!4,colframe=teal!35!black,title={Before you move on}]
What does der Mülleimer mean? (\textquotedblleft{}A bin\textquotedblright{}.) Say it once more.
\end{tcolorbox}
\section[der Besen]{der Besen — a broom}

\label{lesson:GE-C115-besen}



\begin{quote}
Before the new one: say the German for a comb, then the German for a bin.
\end{quote}

\subsection*{You'll want to know: der Besen}
\textbf{der Besen} — \textquotedblleft{}a broom\textquotedblright{}.

\begin{cousinweb}
\textbf{Besen} $\leftarrow$ Old High German \emph{besamo} $\leftarrow$ Proto-Germanic \emph{*besamô}, \textquotedblleft{}broom\textquotedblright{} — inherited, the same word as English \emph{besom}, a broom of twigs. (EtymDB 2.1)
\end{cousinweb}

\begin{grammarlens}[title={What you've built}]
One more word for this chapter.
\end{grammarlens}

\subsection*{Your turn}
\begin{itemize}
\raggedright
  \item \emph{Say it:} \emph{der Besen}
  \item \emph{Say it:} \emph{der Besen}, once more
  \item \emph{Say it:} say \emph{der Besen}
  \item \emph{Recall:} say the German for a comb, then the German for a bin, then say \emph{der Besen} again
\end{itemize}

\begin{tcolorbox}[breakable,colback=teal!4,colframe=teal!35!black,title={Before you move on}]
What does der Besen mean? (\textquotedblleft{}A broom\textquotedblright{}.) Say it once more.
\end{tcolorbox}
\section[die Waschmaschine]{die Waschmaschine — a washing machine}

\label{lesson:GE-C115-waschmaschine}



\begin{quote}
Before the new one: say the German for a bin, then the German for a broom.
\end{quote}

\subsection*{You'll want to know: die Waschmaschine}
\textbf{die Waschmaschine} — \textquotedblleft{}a washing machine\textquotedblright{}.

\begin{cousinweb}
\textbf{Waschmaschine} is the verb \textquotedblleft{}to wash\textquotedblright{}, Old High German \emph{wascan}, English \emph{wash}, + the word for \textquotedblleft{}machine\textquotedblright{} $\leftarrow$ French \emph{machine} $\leftarrow$ Latin \emph{māchina} (EtymDB 2.1).
\end{cousinweb}

\begin{grammarlens}[title={What you've built}]
One more word for this chapter.
\end{grammarlens}

\subsection*{Your turn}
\begin{itemize}
\raggedright
  \item \emph{Say it:} \emph{die Waschmaschine}
  \item \emph{Say it:} \emph{die Waschmaschine}, once more
  \item \emph{Say it:} say \emph{die Waschmaschine}
  \item \emph{Recall:} say the German for a bin, then the German for a broom, then say \emph{die Waschmaschine} again
\end{itemize}

\begin{tcolorbox}[breakable,colback=teal!4,colframe=teal!35!black,title={Before you move on}]
What does die Waschmaschine mean? (\textquotedblleft{}A washing machine\textquotedblright{}.) Say it once more.
\end{tcolorbox}
\section[der Sessel]{der Sessel — an armchair}

\label{lesson:GE-C115-sessel}



\begin{quote}
Before the new one: say the German for a broom, then the German for a washing machine.
\end{quote}

\subsection*{You'll want to know: der Sessel}
\textbf{der Sessel} — \textquotedblleft{}an armchair\textquotedblright{}.

\begin{cousinweb}
\textbf{Sessel} $\leftarrow$ Middle High German \emph{sezzel} $\leftarrow$ Proto-Germanic \emph{*setlaz}, \textquotedblleft{}seat\textquotedblright{} — inherited, the same word as English \emph{settle}, a long wooden bench. (EtymDB 2.1)
\end{cousinweb}

\begin{grammarlens}[title={What you've built}]
One more word for this chapter.
\end{grammarlens}

\subsection*{Your turn}
\begin{itemize}
\raggedright
  \item \emph{Say it:} \emph{der Sessel}
  \item \emph{Say it:} \emph{der Sessel}, once more
  \item \emph{Say it:} say \emph{der Sessel}
  \item \emph{Recall:} say the German for a broom, then the German for a washing machine, then say \emph{der Sessel} again
\end{itemize}

\begin{tcolorbox}[breakable,colback=teal!4,colframe=teal!35!black,title={Before you move on}]
What does der Sessel mean? (\textquotedblleft{}An armchair\textquotedblright{}.) Say it once more.
\end{tcolorbox}
